\documentclass[12pt,a4paper]{article}
\usepackage[utf8]{inputenc}
\usepackage[spanish]{babel}
\usepackage{graphicx}
\usepackage{hyperref}
\usepackage{geometry}
\usepackage{booktabs}
\usepackage{listings}
\usepackage{xcolor}
\usepackage{amsmath}
\usepackage{algorithm}
\usepackage{algpseudocode}

\geometry{margin=2.5cm}

\lstset{
    basicstyle=\small\ttfamily,
    breaklines=true,
    frame=single,
    backgroundcolor=\color{gray!10}
}

\title{Fase 3: Algoritmos de Ruteo Resiliente\\Sistema de Ruteo ante Inundaciones}
\author{Felipe Aros}
\date{Noviembre 2025}

\begin{document}

\maketitle

\section{Modelo de Probabilidad de Falla}

Cada arista $e$ de la red tiene una probabilidad de falla $p_e$ calculada como:

\begin{equation}
p_e = w_1 \cdot f_{inund} + w_2 \cdot f_{dga} + w_3 \cdot f_{rep} + w_4 \cdot f_{lluvia}
\end{equation}

Donde los pesos son:
\begin{itemize}
    \item $w_1 = 0.4$ (inundacion historica)
    \item $w_2 = 0.3$ (proximidad a estacion DGA)
    \item $w_3 = 0.2$ (reportes ciudadanos)
    \item $w_4 = 0.1$ (precipitacion)
\end{itemize}

Cada factor $f_i$ se calcula con funcion exponencial decreciente:

\begin{equation}
f_i = e^{-d_i / r_i}
\end{equation}

Donde $d_i$ es la distancia a la amenaza y $r_i$ el radio de influencia.

\section{Algoritmo 1: Dijkstra Baseline}

Ruta mas corta usando solo distancia como costo.

\begin{lstlisting}[language=SQL]
SELECT * FROM pgr_dijkstra(
    'SELECT id, source, target, length_m AS cost
     FROM infra_aristas',
    source_node, target_node
);
\end{lstlisting}

\textbf{Endpoint}: \texttt{/api/route/baseline}

\section{Algoritmo 2: Dijkstra Resiliente}

Incorpora probabilidad de falla en el costo:

\begin{equation}
cost_e = length_e \cdot (1 + k \cdot p_e)
\end{equation}

Donde $k$ es el factor de penalizacion (default: 5.0).

\begin{lstlisting}[language=SQL]
SELECT * FROM pgr_dijkstra(
    'SELECT id, source, target,
            length_m * (1 + 5.0 * p_fallo_arista) AS cost
     FROM infra_aristas',
    source_node, target_node
);
\end{lstlisting}

\textbf{Endpoint}: \texttt{/api/route/resilient?k=5.0}

\section{Algoritmo 3: CPLEX MILP}

Formulacion como problema de programacion lineal entera mixta:

\subsection{Variables}
\begin{itemize}
    \item $x_e \in \{0, 1\}$: Variable binaria indicando si la arista $e$ pertenece a la ruta
\end{itemize}

\subsection{Funcion Objetivo}
\begin{equation}
\min \sum_{e \in E} x_e \cdot (length_e + \lambda \cdot length_e \cdot p_e)
\end{equation}

\subsection{Restricciones de Flujo}
\begin{equation}
\sum_{e \in out(v)} x_e - \sum_{e \in in(v)} x_e =
\begin{cases}
1 & \text{si } v = source \\
-1 & \text{si } v = target \\
0 & \text{en otro caso}
\end{cases}
\end{equation}

\subsection{Restricciones Opcionales}
\begin{align}
\sum_{e \in E} x_e \cdot p_e &\leq max\_risk \cdot |E| \\
\sum_{e \in E} x_e \cdot length_e &\leq max\_distance
\end{align}

\begin{lstlisting}[language=Python]
from docplex.mp.model import Model

m = Model(name="resilient_routing")
x = {e['id']: m.binary_var(name=f"x_{e['id']}")
     for e in aristas}

# Funcion objetivo
m.minimize(m.sum(
    x[e['id']] * (e['length_m'] + lambda_risk * e['length_m'] * e['p_fallo'])
    for e in aristas
))

# Restricciones de flujo
for v in nodos:
    flow = sum(x[e] for e in out_edges[v]) - sum(x[e] for e in in_edges[v])
    if v == source:
        m.add_constraint(flow == 1)
    elif v == target:
        m.add_constraint(flow == -1)
    else:
        m.add_constraint(flow == 0)

solution = m.solve()
\end{lstlisting}

\textbf{Endpoint}: \texttt{/api/route/optimize?lambda\_risk=5.0}

\section{Algoritmo 4: A* Metaheuristica}

A* modificado con heuristica que considera riesgo.

\subsection{Funcion de Costo}
\begin{equation}
g(n) = \sum_{e \in path(s,n)} length_e \cdot (1 + w_{risk} \cdot p_e)
\end{equation}

\subsection{Heuristica}
\begin{equation}
h(n) = w_h \cdot dist_{euclidiana}(n, target)
\end{equation}

\begin{algorithm}
\caption{A* Resiliente}
\begin{algorithmic}[1]
\State $open \gets \{source\}$, $closed \gets \emptyset$
\State $g[source] \gets 0$
\State $f[source] \gets h(source)$
\While{$open \neq \emptyset$}
    \State $current \gets \arg\min_{n \in open} f[n]$
    \If{$current = target$}
        \State \Return reconstruir\_ruta()
    \EndIf
    \State $open \gets open \setminus \{current\}$
    \State $closed \gets closed \cup \{current\}$
    \ForAll{$neighbor \in vecinos(current)$}
        \If{$neighbor \in closed$}
            \State continue
        \EndIf
        \State $tentative\_g \gets g[current] + cost(current, neighbor)$
        \If{$tentative\_g < g[neighbor]$}
            \State $g[neighbor] \gets tentative\_g$
            \State $f[neighbor] \gets g[neighbor] + h(neighbor)$
            \State $parent[neighbor] \gets current$
            \State $open \gets open \cup \{neighbor\}$
        \EndIf
    \EndFor
\EndWhile
\end{algorithmic}
\end{algorithm}

\textbf{Endpoint}: \texttt{/api/route/metaheuristic?risk\_weight=3.0}

\section{Simulacion de Fallas}

La simulacion activa fallas aleatoriamente segun la probabilidad de cada elemento:

\begin{lstlisting}[language=SQL]
INSERT INTO sim_fallas_activas (simulation_id, entity_type, entity_id,
                                 p_fallo, random_value, is_failed)
SELECT
    gen_random_uuid(),
    'arista',
    id,
    p_fallo_arista,
    random(),
    random() < p_fallo_arista
FROM infra_aristas;
\end{lstlisting}

\textbf{Endpoint}: \texttt{POST /api/simulate-failures}

\section{Comparacion de Algoritmos}

El endpoint \texttt{/api/route/compare} ejecuta los 4 algoritmos simultaneamente y retorna:

\begin{table}[h]
\centering
\begin{tabular}{@{}lcccc@{}}
\toprule
\textbf{Algoritmo} & \textbf{Distancia} & \textbf{Riesgo} & \textbf{Tiempo} & \textbf{Optimalidad} \\
\midrule
Baseline & Minima & Alto & Rapido & Local \\
Resiliente & Moderada & Bajo & Rapido & Local \\
CPLEX & Optima & Configurable & Lento & Global \\
A* & Buena & Bajo & Moderado & Heuristica \\
\bottomrule
\end{tabular}
\caption{Comparacion de caracteristicas de los algoritmos}
\end{table}

\section{Interfaz Web}

La interfaz permite:
\begin{itemize}
    \item Seleccionar origen/destino (GPS, geocodificacion, ID manual)
    \item Visualizar las 4 rutas simultaneamente
    \item Activar/desactivar capas de amenazas
    \item Simular fallas con diferentes tasas (10\%, 30\%, 50\%)
    \item Comparar metricas en tabla
\end{itemize}

\end{document}
